% ==========================================
% BAB V RENCANA SELANJUTNYA
% ==========================================
\setcounter{chapter}{4}
\chapter{RENCANA SELANJUTNYA}
\label{chap:rencana-selanjutnya}
Bab ini berisi rencana implementasi yang akan dilakukan untuk merealisasikan desain yang telah 
dirumuskan pada Desain Solusi. Bab ini menyajikan rencana selanjutnya yang mencakup tiga bagian 
utama, yaitu rencana implementasi prototipe, rencana evaluasi terhadap hasil pengujian, serta 
analisis risiko yang mungkin muncul selama proses tersebut.

\section{Rencana Implementasi}

\subsection{Rencana Implementasi Modul \textit{Mission}}

Implementasi modul \textit{mission} dilakukan dalam lingkungan uji terisolasi untuk memastikan 
seluruh aktivitas data \textit{crawling} dapat diamati tanpa risiko menyebarnya malware uji ke 
lingkungan luar. Lingkungan terisolasi ini memungkinkan proses pengamatan berjalan aman tanpa 
mengganggu sistem utama, serta memberikan ruang untuk melakukan \textit{rollback} apabila terjadi 
gangguan selama pengujian. Rencana implementasi modul \textit{mission} disajikan pada 
Tabel~\ref{tbl:rencana-mission}.

\begin{table}[H]
\centering
\begin{tabular}{|p{3.5cm}|p{8cm}|p{3cm}|}
\hline
\textbf{Langkah} & \textbf{Penjelasan} & \textbf{Tools} \\ 
\hline

Penentuan target data \textit{crawling} &
Mengidentifikasi jenis data yang akan diambil, yaitu hasil \textit{keylogging}, metadata proses sistem, artefak sistem tambahan, dan aktivitas \textit{cache} internet \textit{browser}. &
Data \textit{dummy} \\
\hline

Implementasi \textit{fileless execution} &
Menjalankan modul \textit{mission} sepenuhnya di memori dengan \textit{PowerShell}/\textit{WMI} agar tidak meninggalkan artefak \textit{disk}. &
\textit{PowerShell}, \textit{WMI} \\
\hline

Pemanfaatan \textit{LOLBins} untuk proses \textit{crawling} &
Memanfaatkan perintah bawaan Windows dan query \textit{WMI} untuk membaca direktori, \textit{registry} ringan, serta metadata sistem. Pendekatan ini menyamarkan \textit{crawling} sebagai aktivitas legal. &
Windows native tools \\
\hline

Pengaturan pola \textit{crawling low and slow} &
Mengatur intensitas \textit{crawling} agar berjalan perlahan dan acak, dengan interval yang tidak konstan serta pembacaan \textit{directory} per-\textit{batch} kecil. Tujuannya meniru pola pengguna normal dan memperkecil risiko \textit{alert} dari sistem keamanan. &
Scheduler \\
\hline

Pembentukan \textit{buffer} data in memory &
Data hasil \textit{crawling} disimpan sementara ke dalam \textit{buffer} di memori sebelum dikirimkan ke modul eksfiltrasi. \textit{Buffer} dibersihkan secara berkala untuk menghindari jejak forensik. &
Memory buffer module \\
\hline

\textit{Obfuscation} ringan pada instruksi \textit{crawling} &
\textit{Payload} dan instruksi \textit{crawling} diberi \textit{encoding} atau variasi nama parameter untuk mengurangi deteksi \textit{signature}, tanpa menimbulkan pola anomali baru. &
\textit{Encoder} \\
\hline

Verifikasi fungsi modul \textit{mission} &
Menguji apakah \textit{crawling} berjalan stabil, tidak memicu \textit{alert} awal, dan \textit{buffer} terbentuk dengan benar dalam lingkungan terisolasi. Dilakukan sebelum integrasi dengan modul eksfiltrasi. &
Host Windows 10/11 \\
\hline

\end{tabular}
\captionsetup{justification=raggedright, singlelinecheck=false}
\caption{Rencana Implementasi Modul \textit{Mission}}
\label{tbl:rencana-mission}
\end{table}


\subsection{Rencana Implementasi Modul \textit{Exfiltration}}

Modul \textit{exfiltration} bertugas mengirimkan data hasil proses mission ke server penerima. 
Rencana implementasi disusun untuk memastikan proses pengiriman data berjalan terstruktur, aman, dan 
konsisten. Rencana implementasi modul \textit{exfiltration} disajikan pada Tabel~\ref{tbl:rencana-exfiltration}.

\begin{longtable}{|p{3.5cm}|p{6cm}|p{3cm}|}
\captionsetup{justification=raggedright, singlelinecheck=false}
\caption{Rencana Implementasi Modul \textit{Exfiltration}} \label{tbl:rencana-exfiltration} \\
\hline
\textbf{Langkah} & \textbf{Penjelasan} & \textbf{Tools} \\
\hline
\endfirsthead

\hline
\textbf{Langkah} & \textbf{Penjelasan} & \textbf{Tools} \\
\hline
\endhead

\hline
\endfoot

Menentukan server penerima &
Menyusun layanan sederhana pada sisi penerima untuk menerima data hasil eksfiltrasi dari host Windows. &
Python \textit{FastAPI} \\
\hline

Mempersiapkan struktur data yang akan dikirim &
Data hasil \textit{crawling} yang sudah disimpan sementara dalam \textit{buffer in memory} dikemas menjadi format JSON agar mudah dikirim dan tidak menimbulkan \textit{burst traffic} mencolok. &
\textit{Modul} formatting JSON \\
\hline

Pengaturan interval pengiriman acak (\textit{low and slow}) &
Setiap \textit{payload} dikirim dalam interval waktu acak dan tidak beraturan sehingga pola komunikasi tidak menunjukkan ritme repetitif yang dapat memicu deteksi anomali. &
Timer/Scheduler \\
\hline

Mengirimkan data melalui kanal komunikasi umum (HTTP/S) &
Eksfiltrasi dilakukan menggunakan \textit{request} HTTP/HTTPS lokal sehingga pola komunikasinya menyerupai aplikasi biasa. &
PowerShell \\
\hline

Pembersihan \textit{buffer} setelah pengiriman &
Setelah setiap potongan data dikirim, \textit{buffer} dihapus dari memori untuk meningkatkan \textit{stealth} dan mengurangi jejak aktivitas pada sistem. &
Memory buffer module \\
\hline

Verifikasi pengiriman data &
Memastikan setiap data berhasil diterima oleh server lokal tanpa \textit{error}. &
Log server \\
\hline

\end{longtable}

\section{Rencana Evaluasi}

\subsection{Parameter dan Metrik Evaluasi}

Parameter dan metrik evaluasi digunakan untuk menentukan aspek apa saja yang perlu diamati selama 
pengujian serta bagaimana kualitas perilaku modul diukur. Penetapan parameter membantu memastikan 
setiap aktivitas \textit{mission} dan \textit{exfiltration} dapat dievaluasi secara objektif, 
sementara metrik yang ditentukan berfungsi sebagai indikator keberhasilan atau kegagalan pada setiap 
skenario uji. 

\begin{table}[H]
\centering
\captionsetup{justification=raggedright,singlelinecheck=false}
\caption{Parameter dan Metrik Evaluasi}
\label{tbl:param-metrik}
\begin{tabular}{|p{3.2cm}|p{6cm}|p{5cm}|}
\hline
\textbf{Parameter Evaluasi} & \textbf{Penjelasan} & \textbf{Metrik yang Diukur} \\
\hline

Deteksi pada tingkat proses 
& 
Mengamati apakah aktivitas modul \textit{mission} terdeteksi oleh produk keamanan.
&
\begin{itemize}[leftmargin=*]
  \item Jumlah \textit{alert}
  \item Jenis \textit{alert} (\textit{warning}/\textit{block})
  \item Kategori deteksi (\textit{behavioral}/\textit{signature based})
\end{itemize}
\\
\hline

Deteksi aktivitas \textit{in-memory}/\textit{fileless}
&
Mengevaluasi apakah eksekusi di memori (tanpa menulis \textit{file}) teridentifikasi oleh sistem.
&
\begin{itemize}[leftmargin=*]
  \item \textit{Event log} ID yang muncul
  \item \textit{Flag} atau indikator dari \textit{anti-spyware}
  \item Anomali proses \textit{in-memory}
\end{itemize}
\\
\hline

Deteksi pada aktivitas \textit{crawling}
&
Melihat apakah enumerasi direktori, pembacaan metadata, atau akses \textit{cache} memicu deteksi.
&
\begin{itemize}[leftmargin=*]
  \item \textit{Alert} atau \textit{warning} saat enumerasi \textit{folder}
  \item Peringatan \textit{suspicious behavior}
  \item Jumlah \textit{event} yang ditandai
\end{itemize}
\\
\hline

Deteksi trafik eksfiltrasi
&
Mengamati apakah \textit{batch} JSON yang dikirim melalui HTTPS dianggap anomali oleh \textit{anti-spyware}.
&
\begin{itemize}[leftmargin=*]
  \item \textit{Alert} jaringan
  \item \textit{Block}/allow status
  \item Jumlah \textit{request} yang ditandai
\end{itemize}
\\
\hline

Keberhasilan eksfiltrasi
&
Mengukur seberapa banyak \textit{batch} yang berhasil dikirim dan diterima oleh server.
&
\begin{itemize}[leftmargin=*]
  \item \textit{Success rate} (\%)
  \item Jumlah \textit{batch} gagal
  \item \textit{Latency} per \textit{batch}
\end{itemize}
\\
\hline

Stabilitas modul selama eksekusi
&
Menilai apakah modul berjalan stabil tanpa mengalami \textit{crash} atau \textit{error} selama skenario uji.
&
\begin{itemize}[leftmargin=*]
  \item \textit{Crash count}
  \item \textit{Error log}
  \item Durasi eksekusi tanpa gangguan
\end{itemize}
\\
\hline

\end{tabular}
\end{table}

\subsection{Kriteria Keberhasilan}

Kriteria keberhasilan ditetapkan untuk menentukan apakah modul \textit{mission} dan modul \textit{exfiltration} 
dapat beroperasi sesuai tujuan penelitian serta mengevaluasi respon produk keamanan terhadap aktivitas 
tersebut. Pengujian dinyatakan berhasil apabila memenuhi kriteria keberhasilan pada Tabel \ref{tbl:kriteria-keberhasilan}.

\begin{table}[H]
\centering
\captionsetup{justification=raggedright,singlelinecheck=false}
\caption{Kriteria Keberhasilan}
\label{tbl:kriteria-keberhasilan}
\begin{tabular}{|p{4.5cm}|p{10.2cm}|}
\hline
\textbf{Aspek yang Dievaluasi} & \textbf{Kriteria Keberhasilan} \\
\hline

Deteksi pada tingkat proses &
\begin{itemize}[leftmargin=*]
    \item Jumlah \textit{alert} proses $\leq 5$ 
    \item Tidak ada \textit{block} total oleh \textit{anti-spyware}
\end{itemize}
\\
\hline

Deteksi aktivitas \textit{in-memory/fileless} &
\begin{itemize}[leftmargin=*]
    \item Tidak muncul \textit{alert} \textit{High Severity}
    \item \textit{Alert} \textit{medium/low} $\leq 2$ kali
\end{itemize}
\\
\hline

Deteksi pada aktivitas \textit{crawling} &
\begin{itemize}[leftmargin=*]
    \item \textit{Alert crawling} $\leq 2$ kali
    \item Tidak ada pemblokiran permanen pada akses direktori
\end{itemize}
\\
\hline

Deteksi trafik eksfiltrasi &
\begin{itemize}[leftmargin=*]
    \item \textit{Alert} jaringan $\leq 2$ kali
    \item Tidak terjadi \textit{connection block}
    \item Minimal 70\% \textit{request} HTTPS sukses
\end{itemize}
\\
\hline

Keberhasilan eksfiltrasi &
\begin{itemize}[leftmargin=*]
    \item \textit{Success rate} $\geq 80\%$
    \item \textit{Batch} gagal $\leq 20\%$
    \item Latensi per \textit{batch} masih dalam rentang wajar untuk jaringan lokal ($\leq 5$ detik)
\end{itemize}
\\
\hline

Stabilitas modul selama eksekusi &
\begin{itemize}[leftmargin=*]
    \item Tidak ada \textit{crash} fatal
    \item Modul tetap berjalan sampai skenario uji selesai
\end{itemize}
\\
\hline

\end{tabular}
\end{table}

\section{Analisis Risiko}

Analisis risiko diperlukan untuk mengidentifikasi potensi masalah yang dapat muncul selama implementasi 
dan pengujian modul \textit{mission} serta modul \textit{exfiltration}. Langkah mitigasi disusun agar 
setiap risiko dapat diminimalkan.

\begin{table}[H]
\centering
\captionsetup{justification=raggedright,singlelinecheck=false}
\caption{Analisis Risiko dan Mitigasi}
\label{tbl:analisis-risiko-mitigasi}
\begin{tabular}{|p{2.7cm}|p{3.8cm}|p{3.3cm}|p{4.2cm}|}
\hline
\textbf{Risiko} & \textbf{Penyebab} & \textbf{Dampak} & \textbf{Mitigasi} \\
\hline

Kegagalan eksfiltrasi
&
Server lokal lambat merespon, koneksi terhambat \textit{firewall}, atau interval pengiriman terlalu rapat
&
Data \textit{batch} hilang atau tidak tercatat
&
Menambah interval acak, memvalidasi \textit{endpoint} sebelum uji, dan menyesuaikan ukuran \textit{batch}
\\
\hline

\textit{Crash} atau \textit{hang} pada modul
&
Modul \textit{mission} atau eksfiltrasi mengalami \textit{error runtime}
&
Skenario uji tidak dapat diselesaikan
&
Melakukan uji coba awal, \textit{logging} sederhana, dan pemantauan penggunaan memori selama eksekusi
\\
\hline

Ketidaksesuaian data
&
Kesalahan saat \textit{formatting} JSON atau data \textit{buffer} tidak bersih
&
\textit{Batch} tidak diterima server atau \textit{parsing} gagal
&
Validasi JSON sebelum pengiriman, gunakan struktur \textit{payload} minimalis
\\
\hline

Gangguan pada \textit{host} Windows selama pengujian
&
Beban \textit{crawling} atau proses eksfiltrasi memicu \textit{throttling} CPU/memori
&
Latensi meningkat dan hasil uji bias
&
Membatasi intensitas \textit{crawling}, memastikan \textit{host} dalam kondisi \textit{idle} saat uji
\\
\hline

\end{tabular}
\end{table}